\subsection{Pole placement and PI control}

The desired closed-loop polynomial is chosen as
\begin{equation}
  q(z)=z^2-0.5z+0.06=(z-0.3)(z-0.2).
  \label{eq:desired-polynomial}
\end{equation}
For \(T=\SI{0.08}{s}\), these discrete poles correspond to continuous poles
close to \(-15\,\si{s^{-1}}\) and \(-20\,\si{s^{-1}}\). The target settling
time is therefore about \(\SI{250}{ms}\).

Writing \(q(z)=z^2+q_1z+q_0\), the controller parameters are computed from the
current plant estimates by
\begin{equation}
  d_0=\frac{q_0+a_0}{b_0},
  \qquad
  d_1=\frac{q_1+1-a_0}{b_0}.
  \label{eq:pole-placement-parameters}
\end{equation}
The implemented controller is the incremental PI law
\begin{equation}
  u(k)=u(k-1)+d_1e(k)+d_0e(k-1),
  \label{eq:incremental-pi}
\end{equation}
where \(e(k)=r(k)-y(k)\). The command is saturated at \(\pm255\). The saturated
value, not the unsaturated internal value, is fed back as \(u(k-1)\) in
\eqref{eq:incremental-pi}; this is the anti-windup mechanism used in the
implementation.

The transient cannot be shaped by the pole locations alone. The closed-loop
transfer function contains the controller zero
\begin{equation}
  z_0=-\frac{d_0}{d_1}.
  \label{eq:controller-zero}
\end{equation}
This is the same structural effect described in the reader as a phase error
caused by the different numerator order. A one-degree-of-freedom controller
does not provide an independent numerator design, and therefore cannot remove
this zero without changing the controller structure.
